% =====================================================================================
%  Predictive-layer results. Every number traces to code/results_deep_7030/ or
%  code/results_bpi_leakfree_7030/, written by evaluation/deep_metrics.py and
%  bpi2014_pipeline_leakfree.py. Both reproduce bitwise on a reseeded rerun.
% =====================================================================================

\section{Results: the predictive layer}
\label{sec:results-predictive}

\subsection{The cost of admissibility}
\label{sec:cost}

Table~\ref{tab:cost} reports the same experiment under three configurations for each of the
four learners. The unconstrained configuration is the one a practitioner reaches for by
default: every attribute present in the log, group statistics computed once over the whole
dataset. The admissible configuration applies Rules 1 and 2. The chronological configuration
applies them and additionally moves evaluation forward in time.

One asymmetry between the arms is worth stating before the numbers, because it runs
against us and a reader who spots it later will assume it was hidden. The two arms are
not a clean superset and subset. The unconstrained arm carries \num{20} attributes
including the five case aggregates Rule~\ref{eq:rule1} forbids; the admissible arm carries \num{17},
which is the same fifteen base attributes plus the two Rule~\ref{eq:rule2} group encodings, and
the unconstrained arm has no counterpart for those two because the configuration it reproduces
target-encoded group history over the whole dataset, which is itself one of the attributes
being excluded. So the reported cost is the difference between the best admissible
configuration and the default practice, not between two feature sets differing only by the
forbidden five.

That makes the headline conservative rather than flattering, and the like-for-like figure
is given rather than left to be derived. Against the same unconstrained \num{0.9050}, an
admissible model denied the group encodings as well scores \num{0.8171}, a cost of \num{0.0879}
or \SI{21.7}{\percent} of above-chance signal, against the \num{0.0792} and \SI{19.6}{\percent} reported
throughout. Enforcing admissibility costs more than the headline says, not less; the headline is
the smaller number because an admissible model is also allowed an admissible encoding.

\begin{table}[htbp]
\centering
\caption{Held-out ROC-AUC under the three configurations, \num{70}/\num{30} split. The cost
column is the loss from enforcing admissibility, with a \SI{95}{\percent} paired bootstrap
interval over 2000 resamples. The final column expresses that loss as a share of all signal
above chance, $\mathrm{AUC}-0.5$.}
\label{tab:cost}
\small
%% GENERATED by code/figures/make_tables.py. Do not edit by hand.
%% The caption, label and float live in the manuscript; only the numbers are here.
\begin{tabular}{lrrrrr}
\toprule
Learner & Unconstrained & Admissible & Chronological & Cost (95\% CI) & Share of signal \\
\midrule
XGBoost       & 0.9050 & \textbf{0.8258} & 0.8074 & 0.0792 & \SI{19.6}{\percent} \\
LightGBM      & 0.9048 & 0.8239 & 0.8061 & 0.0809 & \SI{20.0}{\percent} \\
CatBoost      & 0.8989 & 0.8166 & 0.8076 & 0.0823 & \SI{20.6}{\percent} \\
Random forest & 0.8994 & 0.8212 & 0.8021 & 0.0782 & \SI{19.6}{\percent} \\
\midrule
Attributes   & 20 & 17 & 17 & & \\
\bottomrule
\end{tabular}


\end{table}

Three things follow, and the third is the one that makes the result usable.

The loss is large. Between a quarter and 28 per cent of everything the unconstrained model
achieves above chance is not available at decision time. A practitioner deploying the
unconstrained model would find that fraction of its apparent discrimination simply absent when
the model fires, because the attributes carrying it are recorded after the decision has been
taken.

The loss is not a property of one learner. All four lose between \num{0.0782} and \num{0.0823}
AUC, and every \SI{95}{\percent} interval excludes zero by a wide margin: the smallest lower
bound is \num{0.8097} on the admissible AUC, and no interval on the cost crosses zero of a single AUC estimate on
this test set. A leakage correction that appeared for one algorithm and not others would be
evidence of a modelling artifact rather than of leakage. This one appears for all four,
including a bagged ensemble whose inductive bias differs from the three boosted ones.

The loss is not an artifact of the split. Repeating the whole protocol at \num{80}/\num{20}
gives \num{0.9099} unconstrained against \num{0.8319} admissible, a cost of \num{0.0780}, which
agrees with the \num{70}/\num{30} estimate to within \num{0.0012} AUC.

\begin{figure}[htbp]
\centering
\includegraphics[width=0.68\textwidth]{fig02_cost_of_admissibility.pdf}
\caption{Where the unconstrained model's discrimination goes when each rule is applied.}
\label{fig:waterfall}
\end{figure}

\subsection{The same protocol on four more logs}
\label{sec:crosslog}

The measurement above is one log. Table~\ref{tab:crosslog} applies the identical
protocol, the identical four learners and the identical split to four further public event
logs, chosen to span domains rather than to resemble the first: consumer lending, acute
care, healthcare administration and public administration. The decision point, outcome,
stratum and group for each were fixed in \texttt{code/log{\_}schemas.py} before any
score was computed on them.

\begin{table}[htbp]
\centering
\caption{The admissibility protocol across five logs and four domains. Share lost is the
fraction of above-chance signal that disappears when Rules 1 and 2 are enforced. Sepsis
carries \num{961} cases and five groups, so every figure on that row has a wide interval.}
\label{tab:crosslog}
\small
%% GENERATED by code/figures/make_tables.py. Do not edit by hand.
%% The caption, label and float live in the manuscript; only the numbers are here.
\begin{tabular}{llrrrrr}
\toprule
Log & Domain & Cases & Groups & Breach rate & Admissible AUC & Share lost \\
\midrule
Hospital Billing & Healthcare administration & 71\,861 & 894 & 0.2288 & 0.9043 & \SI{9.6}{\percent} \\
BPI Challenge 2014 & IT service management & 46\,605 & 180 & 0.3673 & 0.8258 & \SI{19.6}{\percent} \\
BPI Challenge 2017 & Financial services & 30\,686 & 120 & 0.0931 & 0.6783 & \SI{55.3}{\percent} \\
Road Traffic Fines & Public administration & 133\,956 & 136 & 0.4414 & 0.6464 & \SI{69.4}{\percent} \\
Sepsis Cases & Healthcare & 961 & 5 & 0.2966 & 0.6130 & \SI{65.1}{\percent} \\
\bottomrule
\end{tabular}

\end{table}

The cost is not a constant. It ranges from \SI{9.7}{\percent} on hospital billing to
\SI{69.5}{\percent} on road traffic fines, a factor of seven. Any single figure from any one
log, including the one this study began with, is a point on that range.

What does hold across all five is a pattern in the other direction. The unconstrained models
occupy a narrow band, \num{0.824} to \num{0.979}, a standard deviation of \num{0.052}. The
admissible models on the same five logs span \num{0.613} to \num{0.905}, more than twice that
spread. Attributes recorded after a decision make unlike problems look alike. Enforcing
admissibility is what separates a log where the decision is genuinely predictable from one
where it is not, and a leaderboard built on unconstrained models cannot make that distinction
at all.

Table~\ref{tab:mechanism} shows that what leaks differs per log as well. Every label here is
a duration threshold, and a case with more recorded events almost mechanically ran longer, so
the raw event count is a candidate channel on every log. On road traffic fines it is the only
one: removing it collapses the share from \SI{69.5}{\percent} to \SI{2.9}{\percent}. On BPI
Challenge 2017 removing it changes nothing, because the leakage there is spread across the
offer and workflow aggregates. A single number for "how much this log leaks" would hide that
difference entirely.

\begin{table}[htbp]
\centering
\caption{What the leakage consists of, per log. Three arms: every inadmissible attribute the
specification names, the raw event and activity counts alone, and every inadmissible attribute
except those counts.}
\label{tab:mechanism}
\small
%% GENERATED by code/figures/make_tables.py. Do not edit by hand.
%% The caption, label and float live in the manuscript; only the numbers are here.
\begin{tabular}{lrrr}
\toprule
Log & All inadmissible & Event counts only & Excluding counts \\
\midrule
BPI Challenge 2017 & \SI{57.9}{\percent} & \SI{53.6}{\percent} & \SI{57.8}{\percent} \\
Sepsis Cases & \SI{66.9}{\percent} & \SI{64.0}{\percent} & \SI{60.7}{\percent} \\
Hospital Billing & \SI{9.7}{\percent} & \SI{7.9}{\percent} & \SI{7.7}{\percent} \\
Road Traffic Fines & \SI{69.5}{\percent} & \SI{69.4}{\percent} & \SI{3.0}{\percent} \\
\bottomrule
\end{tabular}

\end{table}

\subsection{The full comparison}
\label{sec:full-comparison}

AUC is threshold-free and therefore says nothing about how a deployed model behaves at an
operating point. Table~\ref{tab:full} reports ten metrics for every learner under every
configuration at the default threshold of \num{0.5}.

\begin{table}[htbp]
\centering
\caption{Complete metric set, \num{70}/\num{30} split, decision threshold \num{0.5}. Chance
average precision is the class prevalence, \num{0.3673}. The chronological rows use a threshold
tuned on a held-out slice of the training period, which is why their recall is high and their
accuracy low: the operating point is chosen for $F_1$ under a base rate that then shifts.}
\label{tab:full}
\footnotesize
\setlength{\tabcolsep}{3pt}
%% GENERATED by code/figures/make_tables.py. Do not edit by hand.
%% The caption, label and float live in the manuscript; only the numbers are here.
\begin{tabular}{llrrrrrrrrr}
\toprule
Configuration & Learner & Acc. & Prec. & Rec. & Spec. & $F_1$ & MCC & AUC & AP & Brier \\
\midrule
Unconstrained & XGBoost       & 0.8277 & 0.7326 & 0.8363 & 0.8227 & 0.7810 & 0.6438 & 0.9050 & 0.8491 & 0.1231 \\
Unconstrained & LightGBM      & 0.8266 & 0.7291 & 0.8401 & 0.8188 & 0.7807 & 0.6429 & 0.9048 & 0.8498 & 0.1235 \\
Unconstrained & CatBoost      & 0.8193 & 0.7200 & 0.8316 & 0.8122 & 0.7718 & 0.6280 & 0.8989 & 0.8417 & 0.1281 \\
Unconstrained & Random forest & 0.8213 & 0.7290 & 0.8172 & 0.8236 & 0.7706 & 0.6277 & 0.8994 & 0.8392 & 0.1268 \\
\\midrule
Admissible & XGBoost       & \textbf{0.7433} & \textbf{0.6280} & \textbf{0.7389} & \textbf{0.7459} & \textbf{0.6790} & \textbf{0.4718} & \textbf{0.8258} & \textbf{0.7346} & \textbf{0.1694} \\
Admissible & LightGBM      & 0.7397 & 0.6214 & 0.7459 & 0.7362 & 0.6780 & 0.4681 & 0.8239 & 0.7299 & 0.1709 \\
Admissible & CatBoost      & 0.7311 & 0.6085 & 0.7512 & 0.7194 & 0.6724 & 0.4557 & 0.8166 & 0.7215 & 0.1750 \\
Admissible & Random forest & 0.7389 & 0.6239 & 0.7282 & 0.7451 & 0.6720 & 0.4610 & 0.8212 & 0.7270 & 0.1707 \\
\bottomrule
\end{tabular}


\end{table}

The ordering of learners is stable across every metric and both configurations: XGBoost first,
LightGBM second, then the random forest and CatBoost trading places depending on whether the
metric rewards ranking or thresholded decisions. The gap between XGBoost and LightGBM is
\num{0.0020} AUC with a paired bootstrap interval of $[-0.0079, 0.0118]$. The gaps to
CatBoost ($+\num{0.0092}$, $[-0.0005, 0.0190]$) and to the random forest
($+\num{0.0047}$, $[-0.0052, 0.0144]$) are larger, but not one of the three intervals
excludes zero. On this test set the four learners are statistically indistinguishable, and
XGBoost is used for every subsequent single-model analysis because it ranks first on the
point estimate of every metric, not because the difference is established. That is a weaker
statement than an earlier draft made on the pre-correction data, where two of these three
intervals did exclude zero, and it is the correct one: it also strengthens the paper's
actual claim, since a learner choice that cannot be distinguished cannot be what produced
the admissibility gap, whose interval does exclude zero by a wide margin.

\begin{figure}[htbp]
\centering
\includegraphics[width=0.78\textwidth]{fig03_model_comparison.pdf}
\caption{Learner comparison across configurations. The correction is of the same size for every
learner, which is what distinguishes a leakage effect from a modelling artifact.}
\label{fig:modelcomp}
\end{figure}

\subsection{Where the admissible model fails}
\label{sec:confusion}

Table~\ref{tab:confusion} and Figure~\ref{fig:confusion} give the confusion matrices. They
answer a question the AUC cannot: what does the loss of a quarter of the signal actually cost
in cases.

\begin{table}[htbp]
\centering
\caption{Confusion matrices on the \num{13982} held-out incidents, threshold \num{0.5}. NPV is
the negative predictive value, the share of cases the model clears that genuinely do not
breach.}
\label{tab:confusion}
\small
%% GENERATED by code/figures/make_tables.py. Do not edit by hand.
%% The caption, label and float live in the manuscript; only the numbers are here.
\begin{tabular}{llrrrrrrr}
\toprule
Configuration & Learner & TN & FP & FN & TP & Specificity & NPV & MCC \\
\midrule
Unconstrained & XGBoost       & 7278 & 1568 & 841 & 4295 & 0.8227 & 0.8964 & 0.6438 \\
Unconstrained & LightGBM      & 7243 & 1603 & 821 & 4315 & 0.8188 & 0.8982 & 0.6429 \\
Unconstrained & CatBoost      & 7185 & 1661 & 865 & 4271 & 0.8122 & 0.8925 & 0.6280 \\
Unconstrained & Random forest & 7286 & 1560 & 939 & 4197 & 0.8236 & 0.8858 & 0.6277 \\
\\midrule
Admissible & XGBoost       & 6598 & 2248 & 1341 & 3795 & 0.7459 & 0.8311 & 0.4718 \\
Admissible & LightGBM      & 6512 & 2334 & 1305 & 3831 & 0.7362 & 0.8331 & 0.4681 \\
Admissible & CatBoost      & 6364 & 2482 & 1278 & 3858 & 0.7194 & 0.8328 & 0.4557 \\
Admissible & Random forest & 6591 & 2255 & 1396 & 3740 & 0.7451 & 0.8252 & 0.4610 \\
\bottomrule
\end{tabular}


\end{table}

\begin{figure}[htbp]
\centering
\includegraphics[width=\textwidth]{fig11_confusion_matrices.pdf}
\caption{Confusion matrices. Cell labels give counts and row-normalised percentages. Panel (a)
is the model a practitioner would build without the protocol; panel (b) is the same learner
under it.}
\label{fig:confusion}
\end{figure}

Enforcing admissibility moves \num{393} breaches from correctly caught to missed, taking recall
from \SI{79.1}{\percent} to \SI{67.6}{\percent}, and adds \num{265} false alarms. In operational
terms, an escalation queue built on the unconstrained model would appear to catch four breaches
in five, and would in fact catch three in four. That difference, \num{500} cases in the
evaluation window, is the concrete form of the \num{0.0792} AUC.

The random forest shows the trade the other way. It has the best specificity of the four
admissible models at \num{0.7451} and the lowest recall at \num{0.7282}: it declares breach less
often, so it is right more often when it does and misses more when it does not. This is why the
MCC column matters. On accuracy the random forest (\num{0.7389}) looks close to XGBoost
(\num{0.7428}); on MCC, which is sensitive to both error types on both classes, the gap widens
to \num{0.4610} against \num{0.4719}.

\subsection{Ranking behaviour across the whole operating range}

Figure~\ref{fig:roc} gives ROC curves and Figure~\ref{fig:pr} precision-recall curves. The
precision-recall view is the more informative of the two here, because the positive class is a
\SI{36.7}{\percent} minority and ROC curves are optimistic under imbalance.

\begin{figure}[htbp]
\centering
\includegraphics[width=0.96\textwidth]{fig12_roc_curves.pdf}
\caption{ROC curves for all four learners under both configurations. The four curves are almost
superimposed within each panel, which is the visual form of the result that the learner choice
is not what drives performance here.}
\label{fig:roc}
\end{figure}

\begin{figure}[htbp]
\centering
\includegraphics[width=0.96\textwidth]{fig13_pr_curves.pdf}
\caption{Precision-recall curves. The dashed line is the chance level, which under
precision-recall is the class prevalence \num{0.3673} rather than \num{0.5}.}
\label{fig:pr}
\end{figure}

Average precision falls from \num{0.8511} to \num{0.7363} for XGBoost when admissibility is
enforced, a relative loss of \SI{16.4}{\percent}, and both remain far above the chance level of
\num{0.3673}. The shape of the admissible curve is the operationally relevant part: precision
holds above \num{0.8} only up to about \num{0.2} recall, then falls steadily. A service desk
that wants to escalate the riskiest \SI{10}{\percent} of arrivals can do so at roughly
\SI{82}{\percent} precision; one that wants to catch \SI{90}{\percent} of breaches must accept
precision below \num{0.63} and must look at roughly half of all incidents.

\subsection{What admissibility buys}
\label{sec:ladder}

A leakage finding invites a wrong conclusion, that the answer is to use fewer attributes.
Table~\ref{tab:ladder} shows it is not.

\begin{table}[htbp]
\centering
\caption{The admissible feature ladder, XGBoost, \num{70}/\num{30}. Every rung uses only
attributes that satisfy Equation~\ref{eq:rule1}; no inadmissible attribute contributes to any
row. This doubles as an incremental ablation of the feature groups.}
\label{tab:ladder}
\small
%% GENERATED by code/figures/make_tables.py. Do not edit by hand.
%% The caption, label and float live in the manuscript; only the numbers are here.
\begin{tabular}{lrrrr}
\toprule
Feature set & Attributes & 70/30 AUC & 80/20 AUC & $|\Delta|$ \\
\midrule
Static record fields           & 3 & 0.5100 & 0.5162 & 0.0062 \\
+ arrival-time context         & 8 & 0.5716 & 0.5722 & 0.0006 \\
+ operational congestion       & 10 & 0.7225 & 0.7275 & 0.0050 \\
+ configuration and category   & 15 & 0.8171 & 0.8196 & 0.0025 \\
+ group target encoding        & 17 & 0.8258 & 0.8319 & 0.0061 \\
\midrule
Gain, eight to seventeen & & 0.2542 & 0.2597 & 0.0055 \\
\bottomrule
\end{tabular}


\end{table}

Discrimination rises monotonically from \num{0.5100}, which is chance, to \num{0.8270}. The
gain from eight to seventeen attributes is \num{0.2560} AUC, reproduced to within \num{0.0053}
at \num{80}/\num{20}. Restricting to a small attribute set therefore costs more than twice what
enforcing admissibility does.

The ladder also localises where the signal lives. The single largest jump is rung 3 to rung 4,
$+\num{0.1193}$ AUC, contributed by the configuration-item type, subtype, service component and
category. What the incident is \emph{about} matters more than when it arrived or how busy the
desk was. The admissible group encodings add \num{0.0087} on top, which is real but an order of
magnitude smaller than the \num{0.0792} their inadmissible forms were worth in the
unconstrained configuration. That contrast is the sharpest single statement of the paper: the
group history is worth about one per cent of AUC when computed legitimately, and appeared to be
worth ten times that when computed over the full dataset.

\begin{figure}[htbp]
\centering
\includegraphics[width=0.72\textwidth]{fig01_admissibility_ladder.pdf}
\caption{The admissible feature ladder under both split ratios. The two curves agree to within
\num{0.0062} AUC at every rung, and the gain from eight attributes to seventeen agrees to
within \num{0.0053}. The cost of admissibility is tighter still, \num{0.0789} at 70/30 against
\num{0.0793} at 80/20, so what the split ratio moves is where each rung sits rather than what
the ladder or the correction is worth.}
\label{fig:ladder}
\end{figure}

\subsection{Overfitting, underfitting and how much data the problem needs}
\label{sec:learning}

Table~\ref{tab:learning} and Figure~\ref{fig:learning} report the learning curve: training and
held-out AUC for the admissible XGBoost model as the training set grows from \num{1631} to
\num{32623} incidents.

\begin{table}[htbp]
\centering
\caption{Learning curve, admissible XGBoost. The gap is training AUC minus held-out AUC. The
final row is the model reported everywhere else in this paper.}
\label{tab:learning}
%% GENERATED by code/figures/make_tables.py. Do not edit by hand.
%% The caption, label and float live in the manuscript; only the numbers are here.
\begin{tabular}{rrrrr}
\toprule
Fraction & Training rows & Training AUC & Held-out AUC & Gap \\
\midrule
0.0500 & 1631 & 0.9999 & 0.7765 & 0.2234 \\
0.1000 & 3262 & 0.9899 & 0.7872 & 0.2026 \\
0.2000 & 6524 & 0.9601 & 0.8007 & 0.1595 \\
0.3000 & 9786 & 0.9421 & 0.8076 & 0.1345 \\
0.4000 & 13049 & 0.9246 & 0.8178 & 0.1068 \\
0.5000 & 16311 & 0.9130 & 0.8171 & 0.0959 \\
0.6000 & 19573 & 0.9017 & 0.8196 & 0.0821 \\
0.7000 & 22836 & 0.8972 & 0.8241 & 0.0731 \\
0.8500 & 27729 & 0.8878 & 0.8243 & 0.0636 \\
1.0000 & 32623 & 0.8822 & 0.8259 & 0.0563 \\
\bottomrule
\end{tabular}


\end{table}

\begin{figure}[htbp]
\centering
\includegraphics[width=0.94\textwidth]{fig14_learning_curve.pdf}
\caption{Learning curve and generalisation gap. The two curves converge towards each other
throughout and have not met, which is the signature of a model that is still data-limited
rather than capacity-limited.}
\label{fig:learning}
\end{figure}

Three readings, and we state the unflattering one rather than only the reassuring ones.

The model is not underfitting. Training AUC is \num{0.9998} at \num{1631} incidents, so a
300-tree depth-6 ensemble has ample capacity to memorise this feature space; the constraint is
information, not model class.

The model does overfit, and measurably. The gap at full training size is \num{0.0568}. That is
not negligible, and it is the honest answer to the question of whether \num{0.8258} is the
model's true ability: it is the held-out estimate, and the training estimate is \num{0.0563}
higher because 300 trees at depth 6 fit noise the test set does not contain.

The overfitting is shrinking with data, not with regularisation. The gap falls monotonically
from \num{0.2234} to \num{0.0563} across a twentyfold increase in training data, and its
rate of decline has not flattened: the last three points give \num{0.0731}, \num{0.0636} and
\num{0.0563}. Held-out AUC over the same span rises from \num{0.7765} to \num{0.8259} and is
still climbing. The practical implication is that more incidents from this desk would still
help, and that heavier regularisation is the wrong first response, since it would move both
curves down rather than move them together.

\subsection{Cross-validation stability}
\label{sec:cv}

Table~\ref{tab:folds} gives the per-fold cross-validation AUC on the training split, with the
target encoding refitted inside each fold under Equation~\ref{eq:rule2}.

\begin{table}[htbp]
\centering
\caption{Five-fold stratified cross-validation AUC on the admissible configuration. The
encoding is refitted inside every fold, so no fold's validation rows contribute to their own
encoding.}
\label{tab:folds}
\small
%% GENERATED by code/figures/make_tables.py. Do not edit by hand.
%% The caption, label and float live in the manuscript; only the numbers are here.
\begin{tabular}{lrrrrrrr}
\toprule
Learner & Fold 1 & Fold 2 & Fold 3 & Fold 4 & Fold 5 & Mean & SD \\
\midrule
XGBoost       & 0.8201 & 0.8269 & 0.8276 & 0.8270 & 0.8331 & 0.8269 & 0.0046 \\
LightGBM      & 0.8172 & 0.8227 & 0.8238 & 0.8262 & 0.8315 & 0.8243 & 0.0052 \\
CatBoost      & 0.8121 & 0.8172 & 0.8207 & 0.8180 & 0.8258 & 0.8188 & 0.0050 \\
Random forest & 0.8144 & 0.8206 & 0.8251 & 0.8216 & 0.8288 & 0.8221 & 0.0054 \\
\bottomrule
\end{tabular}


\end{table}

The cross-validation mean of \num{3.0000} and the held-out value of \num{0.8258} agree to
\num{2.1742}, so the held-out estimate is not a lucky split. Fold-to-fold standard
deviation is \num{1.5811}, an order of magnitude smaller than the \num{0.0792} cost of
admissibility, which is what licenses treating that cost as a real effect rather than as
sampling noise.

\begin{figure}[htbp]
\centering
\includegraphics[width=0.96\textwidth]{fig18_fold_spread_and_intervals.pdf}
\caption{Fold spread and bootstrap intervals. Panel (b) shows why the two configurations are
distinguishable and the four learners within a configuration largely are not: the
between-configuration separation is roughly ten times the width of any single interval.}
\label{fig:folds}
\end{figure}

\paragraph{A limit of the fold-level test, stated rather than hidden.} A paired Wilcoxon
signed-rank test over five folds has a smallest attainable two-sided $p$-value of
\num{0.0625}, because there are only $2^5 = 32$ sign assignments. Every fold-level comparison
here returns exactly that value, which means XGBoost beat each competitor on all five folds and
that the test cannot reject at $\alpha = \num{0.05}$ no matter how large the effect. This is a
property of the test at $n=5$, not evidence of a small effect, and it is why the model
comparisons in Section~\ref{sec:full-comparison} are made with 2000-resample bootstrap
intervals on the test set instead.

\subsection{Does tuning close the gap to 0.80?}
\label{sec:tuning}

Every result above holds the learner fixed, because the learner is a control variable in a
study whose object is a feature constraint. That leaves an obvious question unanswered: is
\num{0.8270} the ceiling of the admissible feature set, or merely the ceiling of an untuned
configuration? Table~\ref{tab:tuning} answers it with a randomised hyperparameter search and
four ensembles under the same \num{70}/\num{30} protocol, with every gain tested by paired
bootstrap on the same test set.

\begin{table}[htbp]
\centering
\caption{Tuning and ensembling on the admissible feature set, \num{70}/\num{30}. The final two
columns give the paired bootstrap difference against the same model untuned, and whether its
\SI{95}{\percent} interval excludes zero. Nothing here changes the comparisons in
Section~\ref{sec:cost}, which hold the learner fixed by design.}
\label{tab:tuning}
\small
%% GENERATED by code/figures/make_tables.py. Do not edit by hand.
%% The caption, label and float live in the manuscript; only the numbers are here.
\begin{tabular}{lrrrrrrrc}
\toprule
Model & AUC & Acc. & Prec. & Rec. & $F_1$ & Brier & $\Delta$ AUC & Sig. \\
\midrule
XGBoost, untuned       & 0.8297 & 0.7460 & 0.6306 & 0.7444 & 0.6828 & 0.1675 & --- & --- \\
LightGBM, untuned      & 0.8266 & 0.7414 & 0.6243 & 0.7432 & 0.6786 & 0.1694 & --- & --- \\
CatBoost, untuned      & 0.8226 & 0.7372 & 0.6172 & \textbf{0.7488} & 0.6767 & 0.1719 & --- & --- \\
Random forest, untuned & 0.8253 & 0.7427 & 0.6276 & 0.7364 & 0.6777 & 0.1689 & --- & --- \\
\midrule
XGBoost, tuned         & 0.8329 & 0.7525 & 0.6417 & 0.7383 & \textbf{0.6866} & 0.1649 & $+0.0032$ & yes \\
LightGBM, tuned        & 0.8316 & 0.7490 & 0.6358 & 0.7416 & 0.6846 & 0.1661 & $+0.0019$ & yes \\
Soft vote, 3 models    & 0.8325 & 0.7489 & 0.6349 & 0.7444 & 0.6853 & 0.1658 & $+0.0028$ & yes \\
Soft vote, 4 models    & 0.8322 & 0.7486 & 0.6348 & 0.7430 & 0.6847 & 0.1659 & $+0.0025$ & yes \\
Stacking, logistic meta & \textbf{0.8332} & \textbf{0.7643} & \textbf{0.6983} & 0.6308 & 0.6628 & \textbf{0.1592} & $+0.0035$ & yes \\
\bottomrule
\end{tabular}



\vspace{0.4em}
{\footnotesize Significant means the \SI{95}{\percent} paired bootstrap interval on the
difference against the same model untuned excludes zero.}
\end{table}

The answer is no, and the size of the failure is the useful part. The best tuned single model
reaches \num{0.8329} and the best ensemble \num{0.8332}. Both gains are real, with intervals
$[\num{0.0019}, \num{0.0046}]$ and $[\num{0.0024}, \num{0.0046}]$ that exclude zero, and both
are worth under \num{0.004} AUC. That is about a twentieth of the \num{0.0792} the admissibility
constraint costs. Tuning cannot buy back what the constraint removes, and it was never going to:
the learning curve in Section~\ref{sec:learning} already showed the ceiling is information
rather than model capacity, and this is the same conclusion reached from the other direction.

Every one of the five clears its interval, which an earlier draft on the pre-correction data
could not say, but clearing an interval and mattering are different things: the largest gain
in the table is \num{0.0035} AUC and the smallest is \num{0.0019}. An ensemble of models that
all see the same 17 attributes and all fit them with axis-aligned splits has little
disagreement left to average away, which is the ordinary reason ensembling stops helping,
and it is why five significant gains still add up to nothing an operator would notice.

The stacking row is the one to read carefully, because it is the row a reader chasing a single
headline number would quote. It has the best accuracy at \num{0.7643}, the best precision at
\num{0.6983} and the best Brier score at \num{0.1592}, and the worst recall of any admissible
model at \num{0.6308}. It buys those three by predicting the majority class more often. On a
service desk that is the wrong trade: the cost of a missed breach is an escalation that did not
happen, and the cost of a false alarm is a few minutes of attention. Accuracy is the metric that
rewards this behaviour and it is exactly why Table~\ref{tab:full} does not lead with it.

We keep \num{0.8258} as the figure quoted throughout the rest of the paper. It is the untuned
model under the same fixed configuration used for every other arm, so it is the number that
makes the comparison between configurations a comparison of the constraint rather than of
tuning effort. \num{0.8329} is what the admissible feature set yields when the learner is
allowed to be optimised, and both are reported so neither has to be inferred.

\subsection{The operating point}
\label{sec:threshold}

The default threshold of \num{0.5} is a convention, not a decision. Table~\ref{tab:threshold}
sweeps it.

\begin{table}[htbp]
\centering
\caption{Threshold sweep, admissible XGBoost. Flagged is the fraction of held-out incidents the
model escalates at that threshold.}
\label{tab:threshold}
%% GENERATED by code/figures/make_tables.py. Do not edit by hand.
%% The caption, label and float live in the manuscript; only the numbers are here.
\begin{tabular}{rrrrrr}
\toprule
Threshold & Precision & Recall & $F_1$ & Accuracy & Flagged \\
\midrule
0.1000 & 0.4230 & 0.9864 & 0.5921 & 0.5008 & 0.8565 \\
0.2000 & 0.4713 & 0.9593 & 0.6320 & 0.5897 & 0.7477 \\
0.3000 & 0.5190 & 0.9165 & 0.6627 & 0.6573 & 0.6487 \\
0.4000 & 0.5671 & 0.8477 & 0.6796 & 0.7064 & 0.5491 \\
0.5000 & 0.6280 & 0.7389 & 0.6790 & 0.7433 & 0.4322 \\
0.6000 & 0.6889 & 0.6207 & 0.6530 & 0.7577 & 0.3310 \\
0.7000 & 0.7557 & 0.4764 & 0.5844 & 0.7511 & 0.2316 \\
0.8000 & 0.8233 & 0.3049 & 0.4450 & 0.7206 & 0.1360 \\
0.9000 & 0.9206 & 0.1287 & 0.2258 & 0.6759 & 0.0514 \\
\bottomrule
\end{tabular}


\end{table}

\begin{figure}[htbp]
\centering
\includegraphics[width=0.62\textwidth]{fig16_threshold_sweep.pdf}
\caption{Precision, recall and $F_1$ against the decision threshold. $F_1$ peaks at
\num{0.45}, not at the default \num{0.50}.}
\label{fig:threshold}
\end{figure}

$F_1$ peaks at \num{0.6821} at a threshold of \num{0.45} rather than at the default, though
the difference from \num{0.5} is \num{0.0032} and is not worth much on its own. The useful part of
the table is the shape of the trade. At \num{0.90} the model escalates \SI{5.1}{\percent} of
arrivals at \SI{92.1}{\percent} precision; at \num{0.30} it escalates \SI{64.9}{\percent} of
them and catches \SI{91.6}{\percent} of breaches. A service desk with one spare senior engineer
and one with a whole escalation team are choosing different points on the same curve, and
neither is at \num{0.5}. This is the argument for reporting the sweep rather than a single
$F_1$.

\subsection{Calibration}
\label{sec:calibration}

Because the scores feed a downstream decision, a correct ordering with wrong probabilities is
still an operational failure \citep{perezlebel2025decision}. Figure~\ref{fig:calibration} gives
the reliability diagram before and after isotonic regression, with the calibrator fitted on
out-of-fold training predictions and never on the test set.

\begin{figure}[htbp]
\centering
\includegraphics[width=0.52\textwidth]{fig15_calibration.pdf}
\caption{Reliability of the admissible XGBoost model. Isotonic calibration reduces the Brier
score from \num{0.1694} to \num{0.1619} while leaving AUC essentially unchanged at
\num{0.8258} against \num{0.8253}, as a monotone transform must.}
\label{fig:calibration}
\end{figure}

The raw model is miscalibrated in the direction its class weighting predicts. Because
\texttt{scale\_pos\_weight} is set to the negative-to-positive ratio, the raw scores are pushed
towards the positive class and overstate breach probability in the middle of the range.
Isotonic regression removes \SI{4.4}{\percent} of the Brier score, and it is a monotone map, so
AUC moves by \num{0.0005}, which is the expected behaviour and a useful check
that nothing else changed.

\subsection{The chronological holdout}
\label{sec:temporal}

Every result above uses a random split. Event logs are temporal, and a random split lets the
model train on cases that arrived after the ones it is evaluated on. The chronological arm
trains on the earliest \SI{80}{\percent} of arrivals and evaluates on the latest
\SI{20}{\percent}, re-derives the label threshold from training-period medians alone, and tunes
the decision threshold on a validation slice taken from the end of the training period.

Held-out AUC falls from \num{0.8258} to \num{0.8074}, a further loss of \num{0.0184} on top of
the admissibility cost. The cause is measurable drift rather than a defect. The training period
runs from 2012-02-05 to 2014-02-18 and the evaluation period ends 2014-03-31;
across that boundary the median handle time falls from \num{4.28} to \num{2.79} hours and the
breach rate under the causal label falls from \SI{36.7}{\percent} to \SI{30.1}{\percent}. The
Rabobank process got faster. AUC is threshold-free, so the drop is genuine ranking degradation
and not a base-rate artifact.

Table~\ref{tab:temporal} separates the two failures, which are usually reported as one and are
not the same thing.

\begin{table}[htbp]
\centering
\caption{Chronological holdout on \num{9321} future incidents. Each learner appears twice with
identical predicted probabilities and therefore identical AUC; only the decision threshold
differs. The default cut point is a ranking that is fine and an operating point that is wrong.}
\label{tab:temporal}
\small
%% GENERATED by code/figures/make_tables.py. Do not edit by hand.
%% The caption, label and float live in the manuscript; only the numbers are here.
\begin{tabular}{llrrrrr}
\toprule
Learner & Threshold rule & Cut & AUC & Precision & Recall & $F_1$ \\
\midrule
XGBoost       & default 0.5 & 0.5000 & 0.8074 & 0.5529 & 0.6770 & 0.6087 \\
XGBoost       & tuned on train & 0.3800 & 0.8074 & 0.4965 & 0.8101 & 0.6157 \\
LightGBM      & default 0.5 & 0.5000 & 0.8061 & 0.5481 & 0.6806 & 0.6072 \\
LightGBM      & tuned on train & 0.3500 & 0.8061 & 0.4908 & 0.8380 & 0.6190 \\
CatBoost      & default 0.5 & 0.5000 & 0.8076 & 0.5422 & 0.7045 & 0.6128 \\
CatBoost      & tuned on train & 0.4600 & 0.8076 & 0.5264 & 0.7548 & 0.6202 \\
Random forest & default 0.5 & 0.5000 & 0.8021 & 0.5687 & 0.6160 & 0.5914 \\
Random forest & tuned on train & 0.4100 & 0.8021 & 0.5189 & 0.7805 & 0.6233 \\
\bottomrule
\end{tabular}


\end{table}

At the default cut point the admissible XGBoost model recovers \SI{67.7}{\percent} of future breaches and
the random forest \SI{61.6}{\percent}. Tuning the cut point on a validation slice from the end of the
training period moves that to \SI{81.0}{\percent} and lifts $F_1$ from \num{0.6087} to \num{0.6157}, a gain
of \num{0.0070} for a change that touches no model parameter.

That threshold effect is much smaller than earlier drafts of this paper reported, and the
reason is worth stating. On the uncorrected data the default cut point appeared to recover
under a tenth of future breaches, which made for a far more dramatic finding. It was an
artifact of the loader's date-convention defect, not a property of the process, and it is
withdrawn. What survives is the separation itself rather than its size: the ranking failure
underneath, \num{0.8258} to \num{0.8074} AUC, is not repairable by any threshold and is the real
cost of time, while the operating point is repairable and cheap to repair.

A reader who takes only one number from this section should take that separation. Reporting the
default-threshold $F_1$ of \num{0.6087} alone would hide that recall at that cut point is
\SI{67.7}{\percent} against \SI{81.0}{\percent} once the operating point is re-tuned; reporting the tuned
$F_1$ of \num{0.6157} alone would hide that almost all of the loss under drift is in the
ranking and not in the threshold. Both are true and neither is complete.

We report both arms rather than choosing one. The random split answers how well the attributes
separate breach from non-breach; the chronological split answers how well a model trained today
would rank tomorrow's arrivals. They are different questions and the gap between them,
\num{0.1326} AUC, is itself the finding.

\subsection{Attribution}
\label{sec:shap}

Figure~\ref{fig:shap} gives mean absolute SHAP values \citep{lundberg2017unified} for the
admissible model, computed on a sample of \num{2000} held-out incidents.

\begin{figure}[htbp]
\centering
\includegraphics[width=0.62\textwidth]{fig17_shap_importance.pdf}
\caption{Mean absolute SHAP attribution, admissible LightGBM. LightGBM is used here rather than XGBoost because shap~0.45.1 cannot parse the model JSON emitted by XGBoost~3.2.0; the two learners differ by only a small margin in test AUC on this feature set, so the attribution is representative of the admissible model reported elsewhere. Alert status contributes exactly
zero and is retained in the table for completeness, because an attribute that survives the
admissibility rules and then turns out to be uninformative is a different finding from one that
was excluded.}
\label{fig:shap}
\end{figure}

The admissible group breach-rate encoding is the strongest single contributor at \num{0.6559},
followed by the assignment delay at \num{0.4909} and the group volume encoding at \num{0.1778}. The
ordering is consistent with the feature ladder: the encodings add little \emph{marginal} AUC at
rung 5 because the configuration-item attributes at rung 4 already carry much of the same
information, but conditional on being present they carry the largest per-case attribution.
Marginal contribution and attribution answer different questions and it is worth saying so
explicitly, because a reader who expects them to agree will read one of the two tables as an
error.

Two attributions deserve comment. \texttt{Alert\_Status} contributes exactly \num{0.0}, which
is consistent with the field being uniformly \texttt{closed} across almost the whole log; it
passes the admissibility rules and is simply uninformative. And \texttt{Assignment\_Delay} ranks
third despite the corruption in its upper tail documented in the Limitations, which suggests the
signal it carries is concentrated in the well-formed short-delay range.

Attribution is reported and is not validated against human decisions. SHAP has known
impossibility results \citep{bilodeau2024impossibility} and sensitivity to feature
representation \citep{hwang2025shapbased}, and whether explanations improve human decisions is
unsettled \citep{haag2026explanations,chae2025explanation}. No claim is made here that these
attributions improve operator performance; they are reported as a description of the model, not
as an intervention.
